\documentclass[../../../include/open-logic-section]{subfiles}

\begin{document}

\olfileid{sth}{cardinals}{cp}
\olsection{কান্টরের নীতি}

\olref[ordinals][vn]{sec}-এর আলোচনা মনে করুন। সেখানে আমরা
সুক্রমবিন্যস্ত সেট নিয়ে কথা বলেছিলাম এবং বলেছিলাম, সুক্রমবিন্যাসের
প্রতিনিধি হিসেবে কাজ করবে এমন বস্তু পাওয়া ভালো হবে। সেই ভাবনা
থেকে ক্রমসংখ্যার পরিচয় দিয়েছিলাম। পরে
\olref[ordinals][ordtype]{ordtypesworklikeyouwant}-এ দেখিয়েছি,
ক্রমসংখ্যাগুলি প্রত্যাশামতোই কাজ করে, অর্থাৎ:
\[
	\ordtype{A, <} = \ordtype{B, \lessdot}
	\text{ iff } \tuple{A, <} \isomorphic \tuple{B, \lessdot}.
\]
আরও আগে \olref[sfr][siz][equ]{sec}-এর কথাও মনে করুন। সেখানে
অনানুষ্ঠানিকভাবে সেটের ``আকার'' ধারণাটি এনেছিলাম। বিশেষত,
দুটি সেটের মধ্যে !!a{bijection} $f\colon A\to B$ থাকলেই
তাদের সমসংখ্যক বলেছিলাম, অর্থাৎ $\cardeq{A}{B}$। ধারণাটি
সুক্রমবিন্যাসের চেয়ে স্বভাবতই \emph{সরল}: এখানে আমাদের আগ্রহ
শুধু !!{bijection}s-এ; ক্রমসমরূপতার ক্ষেত্রে যেমন লাগে,
!!{bijection}s কোনো ``কাঠামো রক্ষা করে'' কি না, তা এখানে জরুরি নয়।

এ থেকে একটি স্বাভাবিক ভাবনা আসে। সুক্রমবিন্যাসের মাপ নির্ধারণের
জন্য যেমন বিশেষ বস্তু \emph{ক্রমসংখ্যা} এনেছিলাম, আকারের মাপ
নির্ধারণের জন্য তেমনি বিশেষ বস্তু \emph{অঙ্কবাচক সংখ্যা} আনতে
পারি। এই অধ্যায়ের উদ্দেশ্য সেটিই।

অঙ্কবাচক সংখ্যা কী হবে তা বলার আগে, তাদের জন্য একটি নীতি
স্থির করা দরকার। সেট $X$-এর অঙ্কবাচকতা $\card{X}$ লিখলে আমরা চাই:
\[
	\card{A} = \card{B} \text{ iff } \cardeq{A}{B}.
\]
এটিকে \emph{কান্টরের} নীতি বলব, কারণ সম্ভবত কান্টরই প্রথম
এটি স্পষ্টভাবে উপলব্ধি করেছিলেন। (\olref[hp]{sec}-এ
\emph{হিউমের} নীতির সঙ্গে এর সম্পর্ক নিয়ে আরও বলব।)
তাই আমাদের লক্ষ্য, প্রত্যেক $X$-এর জন্য $\card{X}$ এমনভাবে
সংজ্ঞায়িত করা যাতে কান্টরের নীতি মেলে।

\end{document}
